\hwhat{The question is whether agents co-move because a common input hits everyone (a field) or because they read each other (a coupling). In call cycles, a field moves every agent at its next call. A coupling acts only at the recipient's first call that can read the message. The DQ1 context ledger gives a placebo-corrected discontinuity, hop by hop: calls that start just after a message against calls just before it. The analysis also fits input impulse responses and decomposes co-movement on 71 non-holdout units (35 periods). \textbf{Activity co-movement is the schedule}: day start and stop explain about 0.6 of it beyond a null. \textbf{Talk co-movement is a gated coupling}. A peer message raises the chance that the next call is a talk call (45/71 units, none negative), always at hop 1. It accounts for roughly all talk co-movement. HH167's regime split fails. The holdout is not run.}

\hmodel{Gated linear response (models 02, 01). Agent $i$ makes calls $c$ at times $t_c$. A message posted at $t_m$ enters the context at the first call with $t_c>t_m$ (hop 1). The call outcome $Y_c$ (talk, work, idle) responds to inputs $u_{c-k}$ (bookends, human messages, nudges, platform errors). It also responds to peer messages $m_{c-k}$ first read $k$ hops earlier. $s$ is the call start minus the message time. $e_{it}$ is agent-day-demeaned minute activity or talk. $S$ is the equal-time co-movement.}

\hmath{\vspace{-6pt}\begin{gather*}
\mathrm{logit}\,P(Y_c{=}1)=h_i+\textstyle\sum_k g(k)\,u_{c-k}+\sum_k J(k)\,m_{c-k}\\
J_k=E[Y\,|\,s\!\to\!0^+]-E[Y\,|\,s\!\to\!0^-]-\text{placebo},\quad \delta(h)=\textstyle\sum_{k\le h}J_k\\
S=\textstyle\sum_t\big[(\sum_i e_{it})^2-\sum_i e_{it}^2\big]\big/\sum_t\sum_i e_{it}^2
\end{gather*}\vspace{-8pt}}

\hobs{../figures/summary_kernel_split_col.pdf}{Left: excess talk probability of a recipient's call, by hop after a peer message (IVW pooled; 95\% band). The response is zero at the in-flight call (hop 0) and steps up at the read-out call (hop 1). Regime III (13-s chained calls) drops at hop 2, then plateaus. Regime I (scheduled chat calls) keeps rising. Right: share of activity co-movement explained by the schedule edges (gray), and of talk co-movement by the gated coupling (orange; $>1$ is overshoot).}

\htelos{Do not read synchrony in activity as influence. The scheduler's daily start and stop explain about 0.6 of activity co-movement; agents do not cause it. Influence acts at the recipient's next call and turns a work call into a talk call. In the always-on regime it needs the recipient's name: the jump is 0.17 for named recipients and 0.004 for others. To propagate a message, address the agents by name and expect a response within one call (about 15--25 s; minutes if the agent pauses). The nudger moves single agents and leaves the collective mode unchanged. Bookend messages are announcements, not inputs.}

\hbottom{Activity co-movement is a field set by the scheduler; talk spreads by a coupling gated at one read-out call, and agents rarely relay it further.}

\hresults{\scriptsize\setlength{\tabcolsep}{2pt}\begin{tabular}{@{}p{0.36\columnwidth}p{0.45\columnwidth}p{0.16\columnwidth}@{}}
\textbf{Prediction (dated)} & \textbf{Observed} & \textbf{Verdict}\\\hline
P1 synthetic recovery & $J_1$ bias $-2$ to $-17\%$ (III); dead time 2 $\to$ onset hop 3; nulls 3/35 & passed\\
P2 inputs at read-out & in-room human $J_1>0$, hop 1; nudged agents talk at hop 1 (0.14), not later & mixed\\
P3 low-pass, corner 5--40 min & coherence 8--10$\times$ at 15--60 min; corner ill-defined (undershoot) & mixed\\
P4 field: zero relative lag & onsets IQR 9--22 s ($<2$ cycles); \#38 agents halted at pause & supported\\
P5 gated coupling, hop 1 & $J_1>0$ in 45/71, $<0$ in 0; onset hop 1; no decay in regime I & mostly\\
P6 decomposition & activity field 0.55 (III) \checkmark, 0.51 (I, pred.\ $\le0.2$) $\times$; talk coupling $>$ field 57/68 & mostly\\
HH167 regime split & regime-I calls not message-triggered (ratio 1.03--1.07) & rejected\\
NE43 switch-offs & nudge path gone; coupling $\times0.69$; edge field persists & mixed\\
G51 human relay & plain human messages: hop-1 field; relay underpowered & mixed\\
G38 pause gates & 82\% start before any peer talks; 94\% never call after pause & supported\\
\end{tabular}}

\hobsb{../figures/synthetic_col.pdf}{Synthetic swarms at village sampling with real input schedules (5 seeds each). The read-out jump finds planted coupling with or without fields, slow common drives or edge-clustered inputs. It reads zero for fields and drives. It recovers a 2-hop dead time as a later onset. It goes negative for ungated coupling.}

\hcaveats{Regime-I call starts are latency-placed. The jump survives on logged-start recipients (\#24--\#31; 0.030 vs 0.034), but early regime I is unchecked. The rising regime-I kernel beyond hop 2 has no synthetic validation. Per-unit CIs are mildly anti-conservative (synthetic false positives about 9\%). Decomposition shares are model-based and do not add up, because coupling amplifies fields. $f_C>1$ is overshoot at low talk co-movement. Room changes and roster joins confound NE43. The \#51 target and relay tests have about 100 pairs. No content channel exists yet. The holdout is not run.}

\hnext{Add a content gate with a matched-age control. Validate the regime-I kernel. Compare per-agent ``impedance'' with H29/H32 influence. Confirm on the holdout, including NE23 (nudger off and on).}
